% Pista geo-23s-q5 · Geometria
% Procedència: PAU Catalunya 2023 · Convocatòria de setembre · Sèrie 2 · Qüestió 5
\documentclass[11pt,a4paper]{article}
\usepackage{pista}
\pistainfo{Catalunya 2023}{Convocatòria de setembre}{Sèrie 2}

\begin{document}

\section*{\textcolor{blauPista}{Qüestió 5}}

\noindent Siguin $r_{1}$ i $r_{2}$ les rectes definides per $r_{1}: x - 1 = y = -z$ i per $r_{2}: x = y = z$, respectivament.

\subsection*{\textit{a)} \normalfont Calculeu l'equació paramètrica de la recta que talla perpendicularment les rectes $r_{1}$ i $r_{2}$. \hfill\textnormal{[1,75~punts]}}

\pista{Comencem extraient la informació bàsica: identifica un vector director i un punt de cada recta a partir de les equacions contínues que et donen. La recta perpendicular comuna que busques té dues característiques que la determinen: la \emph{direcció} (ha de ser perpendicular alhora a $\vec{v}_{1}$ i a $\vec{v}_{2}$) i \emph{el fet que talla} les dues rectes. Per a la direcció, quina operació entre dos vectors et dóna directament un vector perpendicular als dos? I per trobar el punt, una manera és la següent: parametritza un punt genèric $P_{1}$ sobre $r_{1}$ (amb un paràmetre $t$) i un punt genèric $P_{2}$ sobre $r_{2}$ (amb un paràmetre diferent $s$), i escriu una equació que diu que el vector $\overrightarrow{P_{1} P_{2}}$ sigui perpendicular alhora a $\vec{v}_{1}$ i a $\vec{v}_{2}$ (és a dir, dos productes escalars iguals a zero). Et quedarà un sistema de dues equacions amb dues incògnites; resol-lo i obtindràs els valors de $t$ i $s$ que et donen els punts $P_{1}$ i $P_{2}$. La recta que els uneix és, precisament, la perpendicular comuna que busques.}

\subsection*{\textit{b)} \normalfont Calculeu la distància entre $r_{1}$ i $r_{2}$. \hfill\textnormal{[0,75~punts]}}

\pista{La distància entre les dues rectes és, exactament, la longitud del segment de la perpendicular comuna que va de $P_{1}$ a $P_{2}$, és a dir, el mòdul del vector $\overrightarrow{P_{1} P_{2}}$. Com que els punts $P_{1}$ i $P_{2}$ ja els has trobat a l'apartat anterior, només et cal calcular aquest mòdul.}

\end{document}
